
考虑引理~\ref{lem:7-3-24}
所得多值截面系统的零集
\[
\bigl([0,1]^{|\vec a|}\times
\mathcal M^{\mathrm{main},\mathfrak d^+}_{k+1}
(\beta;\vec P^+)\bigr)
^{\mathfrak s_{\vec a,\mathfrak d^+,k,\beta,\vec P^+}}.
\]
为简化记号，下文把它写作
\[
\bigl([0,1]^{|\vec a|}\times
\mathcal M^{\mathrm{main},\mathfrak d^+}_{k+1}
(\beta;\vec P^+)\bigr)^{\mathfrak s}.
\]
由引理~\ref{lem:7-3-24}，
它具有与圆盘分层相容的三角剖分。

下一步要选取与遗忘映射相容的三角剖分，
并把该零集的一个单纯子复形塌缩，
使前述情形 1 的贡献作为奇异链严格等于零。
模空间作为带仓西结构的空间，有边界分解
\[
\partial\mathcal M^{\mathrm{main}}_{k+1}(\beta)
=
\bigcup_{\beta_1+\beta_2=\beta}
\bigcup_{k_1+k_2=k+1}
\bigcup_{i=1}^{k_2}
\mathcal M^{\mathrm{main}}_{k_1+1}(\beta_1)
{}_{\operatorname{ev}_0}\times_{\operatorname{ev}_i}
\mathcal M^{\mathrm{main}}_{k_2+1}(\beta_2).
\]
还须引入一些记号。
把
\[
\vec a=(\vec a(1),\ldots,\vec a(|\vec a|))
\in\{1,\ldots,k\}^{|\vec a|}
\]
分解为不交并
$\vec a=\vec a^1\sqcup\vec a^2$。
这里数据 $\vec a,i,k_1,k_2$
按下式确定 $\vec a^1,\vec a^2,m$：
\begin{equation}
\vec a(j)=
\begin{cases}
\vec a^1(j),
&j\le m,\\
\vec a^2(j-m)+i-1,
&m+1\le j\le m+|\vec a^2|,\\
\vec a^1(j-|\vec a^2|+m)+k_1-1,
&m+1+|\vec a^2|+1\le j.
\end{cases}
\tag{7.3.26}\label{eq:7-3-26}
\end{equation}
对每个 $\vec a$，
取 $\mathcal M^{\mathrm{main}}_{k+1}(\beta)$
的一个副本，记作
$\mathcal M^{\mathrm{main},\vec a}_{k+1}(\beta)$。\translatornote{源书第三支的下界
$m+1+|\vec a^2|+1\le j$
遗漏紧接第二支的一个指标；
自变量 $j-|\vec a^2|+m$ 也有疑点。
此处暂照原样保留，见疑误 A-0547。}
考虑交换图：

\par\medskip\noindent\begin{minipage}{\textwidth}
\centering
\resizebox{\textwidth}{!}{$\displaystyle
\begin{array}{ccc}
\mathcal M^{\mathrm{main},\vec a^1}_{k_1+1}(\beta_1)
{}_{\operatorname{ev}_0}\times_{\operatorname{ev}_i}
\mathcal M^{\mathrm{main},\vec a^2}_{k_2+1}(\beta_2)
&\longrightarrow&
\mathcal M^{\mathrm{main},\vec a}_{k+1}(\beta)\\[0.4em]
{\scriptstyle\operatorname{forget}\times\operatorname{forget}}\Big\downarrow
&&
\Big\downarrow{\scriptstyle\operatorname{forget}}\\[0.4em]
\mathcal M^{\mathrm{main}}_{k_1+1-|\vec a^1|}(\beta_1)
{}_{\operatorname{ev}_0}\times_{\operatorname{ev}_i}
\mathcal M^{\mathrm{main}}_{k_2+1-|\vec a^2|}(\beta_2)
&\longrightarrow&
\mathcal M^{\mathrm{main}}_{k+1-|\vec a|}(\beta)
\end{array}$}
\par\smallskip
{\renewcommand{\thefigure}{7.3.1}\refstepcounter{figure}
\normalfont\bfseries 图式~\thefigure\label{diag:7-3-1}}
\end{minipage}\par\medskip

记
\[
[0,1]^{k-\vec a}
=\bigl\{(t_1,\ldots,t_k)\mid
t_{\vec a(i)}=1,\ i=1,\ldots,|\vec a|\bigr\}
\cong[0,1]^{k-|\vec a|}.
\]
把 $\vec a=\vec a^1\sqcup\vec a^2$
按上式分解后，
遗忘映射与例~\ref{ex:7-3-12} 紧前定义的
$I^1_{\vec a^1,\vec a^2}$
诱导映射
\[
\operatorname{forget}\colon
[0,1]^{k-\vec a}\times
\mathcal M^{\mathrm{main},\vec a}_{k+1}
(\beta;\vec P^+)
\longrightarrow
[0,1]^{k-|\vec a|}\times
\mathcal M^{\mathrm{main}}_{k+1-|\vec a|}
(\beta;\vec P).
\]
由相容条件~\ref{cond:7-3-13}，
它与多值截面系统相容。
于是图式~\ref{diag:7-3-1}
和这一相容性给出下一个交换图。
为使图式易读，先列出它的四角对象：
\[
\begin{aligned}
Z_j^+
&=
\bigl([0,1]^{k_j-\vec a^j}\times
\mathcal M^{\mathrm{main},\vec a^j}_{k_j+1}
(\beta_j;\vec P_j^+)\bigr)^{\mathfrak s},
&&j=1,2,\\
Z^+
&=
\bigl([0,1]^{k-\vec a}\times
\mathcal M^{\mathrm{main},\vec a}_{k+1}
(\beta;\vec P^+)\bigr)^{\mathfrak s},\\
Z_j
&=
\bigl([0,1]^{k_j-|\vec a^j|}\times
\mathcal M^{\mathrm{main}}_{k_j+1-|\vec a^j|}
(\beta_j;\vec P_j)\bigr)^{\mathfrak s},
&&j=1,2,\\
Z
&=
\bigl([0,1]^{k-|\vec a|}\times
\mathcal M^{\mathrm{main}}_{k+1-|\vec a|}
(\beta;\vec P)\bigr)^{\mathfrak s}.
\end{aligned}
\]
这里图式中的 $\times_L$ 表示
$\,{}_{\operatorname{ev}_0}\times_{\operatorname{ev}_i}$。

\par\medskip\noindent\begin{minipage}{\textwidth}
\centering
$\displaystyle
\begin{array}{ccc}
Z_1^+\times_L Z_2^+&\longrightarrow&Z^+\\[0.4em]
{\scriptstyle\operatorname{forget}\times\operatorname{forget}}\Big\downarrow
&&
\Big\downarrow{\scriptstyle\operatorname{forget}}\\[0.4em]
Z_1\times_L Z_2&\longrightarrow&Z
\end{array}$
\par\smallskip
{\renewcommand{\thefigure}{7.3.2}\refstepcounter{figure}
\normalfont\bfseries 图式~\thefigure\label{diag:7-3-2}}
\end{minipage}\par\medskip

\begin{lemma}{7.3.27}\label{lem:7-3-27}
可以为扰动模空间
\[
\bigl([0,1]^{|\vec a|}\times
\mathcal M^{\mathrm{main},\mathfrak d^+}_{k+1}
(\beta;\vec P^+)\bigr)^{\mathfrak s}
\]
选取一组三角剖分，
使图式~\ref{diag:7-3-2} 中的箭头\translatornote{源书图式左下角第二因子
仍写 $\vec P_2^+$；
按该角的遗忘映射及右下角输入
$\vec P$，这里写成 $\vec P_2$，
见疑误 A-0546。}
均为分段线性。
还可使这一三角剖分
与相容条件~\ref{cond:7-3-17}
给出的识别相容。
\end{lemma}

\begin{proof}
先看图式~\ref{diag:7-3-2} 的水平箭头。
由圆盘分层的定义容易看出，
其限制逐层是局部微分同胚。
再由引理~\ref{lem:7-3-8} 可知，
垂直箭头在每一层上的限制是浸没。
因此可逐层归纳构造三角剖分。
\end{proof}
